\section*{Chapitre \ref{ch:dv:sabr-and-smile-dynamics} --- SABR et dynamique du smile}

\begin{solution}{exo:dv:sabr-and-smile-dynamics:1}
$\beta=1$, $\nu=0$~: un forward lognormal à volatilité constante $\alpha$, un smile de Black plat à $\alpha$. $\beta=0$,
$\nu=0$~: le modèle de Bachelier, un smile normal plat à $\alpha$, qui, en volatilité de Black, est un skew décroissant à peu près comme
$\alpha/\sqrt{FK}$.
\end{solution}

\begin{solution}{exo:dv:sabr-and-smile-dynamics:2}
$\alpha=0.2\times0.03^{1-\beta}$ et $\sigma_{\mathrm{ATM}}(0.025)=\alpha\times0.025^{\beta-1}$~: 24.0~\% pour $\beta=0$,
21.9~\% pour $\beta=0.5$, 20.0~\% pour $\beta=1$.
\end{solution}

\begin{solution}{exo:dv:sabr-and-smile-dynamics:3}
Chaque échéance est calibrée avec ses propres paramètres, si bien que le modèle n'a jamais à décrire l'évolution de la volatilité d'une échéance à l'autre~:
le retour à la moyenne, qui façonne une structure par terme, n'intervient pas.
\end{solution}

\begin{solution}{exo:dv:sabr-and-smile-dynamics:4}
$dF=\alpha F^\beta dW^1$ et $d\alpha=\nu\alpha dW^2$ avec la corrélation $\rho$~: $\Cov(d\alpha,dF)=\rho\nu\alpha^2F^\beta dt$
et $\Var(dF)=\alpha^2F^{2\beta}dt$, si bien que la pente vaut $\rho\nu/F^\beta$. La sensibilité de l'option à $\alpha$, multipliée par cette pente,
est la part de son risque en $\alpha$ couvrable avec le forward~: $\mathcal V\,\partial_\alpha\sigma_B\,\rho\nu/F^\beta$.
\end{solution}

\begin{solution}{exo:dv:sabr-and-smile-dynamics:5}
À partir de la co-variation historique de la volatilité à la monnaie et du forward~: on régresse la variation de $\ln\sigma_{\mathrm{ATM}}$
sur la variation de $\ln F$~; la pente estime $\beta-1$. Un $\beta$ erroné est compensé par $\rho$ dans l'ajustement, si bien que les
prix ne changent pas, mais les deltas sont faux de la différence de pente du \omterm{def:dv:sabr-and-smile-dynamics:backbone}{backbone} multipliée par le \omterm{def:dv:greeks-and-the-hedging-pnl:greeks}{vega}.
\end{solution}

\begin{solution}{exo:dv:sabr-and-smile-dynamics:6}
70.8, 59.9 et 55.0 points de base par an, contre 28.8~\%, 20.0~\% et 15.8~\% en volatilité de Black. Une volatilité normale vaut
à peu près la volatilité de Black multipliée par $\sqrt{FK}$~; une grande part de la baisse du smile de Black avec le strike
provient du facteur d'échelle $1/\sqrt{FK}$ propre au modèle lognormal, que la cotation normale supprime.
\end{solution}

\begin{solution}{exo:dv:sabr-and-smile-dynamics:7}
Les \omterm{def:dv:sabr-and-smile-dynamics:mvdelta}{deltas de variance minimale} et de Black donnent presque la même erreur (rapport de variance 0.999)~; celui de Hagan est
plus mauvais de 10~\%. Avec une corrélation positive, le \omterm{def:dv:sabr-and-smile-dynamics:backbone}{backbone} et la co-variation de la volatilité se compensent presque exactement à la monnaie.
\end{solution}

\begin{solution}{exo:dv:sabr-and-smile-dynamics:8}
L'ajustement porte sur les strikes et les échéances cotés~; un strike de 1~\% à dix ans est en dehors de ces points, là où la \omterm{def:dv:sabr-and-smile-dynamics:hagan}{formule
de Hagan} peut produire une densité négative (inférieure à environ 1.6~\% dans l'exemple à dix ans du chapitre), si bien que son prix n'est même pas sans arbitrage à cet endroit.
Il faut utiliser la densité sans arbitrage ou un modèle décalé, et prévoir une réserve pour l'extrapolation.
\end{solution}

\begin{solution}{pb:dv:sabr-and-smile-dynamics:1}
\textbf{1.} $\alpha=0.0346$.
\textbf{2.} 28.8~\%, 20.0~\%, 15.8~\%.
\textbf{3.} 59.9 points de base par an.
\textbf{4.} 18.5~\%.
\textbf{5.} Avec le forward~: son minimum et sa forme se déplacent vers la droite, tandis que son niveau glisse le long du \omterm{def:dv:sabr-and-smile-dynamics:backbone}{backbone}.
\textbf{6.} 0.540, 0.580 et 0.461.
\textbf{7.} Le delta de Hagan ajoute le \omterm{def:dv:greeks-and-the-hedging-pnl:greeks}{vega} multiplié par la pente de la volatilité implicite en $F$, positive ici parce que le
point à la monnaie remonte le côté droit du smile lorsque $F$ augmente (la volatilité d'un strike fixe croît avec $F$).
\textbf{8.} Avec $\rho<0$, une hausse de $F$ s'accompagne d'une baisse de $\alpha$, ce qui diminue la valeur du call~; la couverture
doit l'anticiper.
\textbf{9.} $\rho\nu/F^\beta=-0.6\times0.5/\sqrt{0.03}=-1.73$.
\textbf{10.} La part de $d\alpha$ non corrélée à $dF$~: risque de \omterm{def:dv:greeks-and-the-hedging-pnl:greeks}{vega}, couvert avec une autre option.
\textbf{11.} $6.67\times10^{-5}$ (Black), $7.46\times10^{-5}$ (Hagan), $5.99\times10^{-5}$ (variance minimale).
\textbf{12.} 1.25 et 0.81 de celui de Black.
\textbf{13.} Il déplace la volatilité le long du smile à $\alpha$ fixé, alors que dans le modèle $\alpha$ tend à baisser lorsque
$F$ monte~: il corrige dans le mauvais sens.
\textbf{14.} Les deltas de Hagan et de variance minimale coïncident, avec 1~\% de variance en moins que celui de Black dans cet exemple.
\textbf{15.} Par $1-\sqrt{0.81}$, soit environ 10~\%.
\textbf{16.} Il faut présenter le delta du modèle utilisé pour la couverture, et traiter le \omterm{def:dv:sabr-and-smile-dynamics:mvdelta}{delta de variance minimale} si la dynamique du modèle
est jugée fiable~; les autres se présentent comme des sensibilités à la dynamique.
\textbf{17.} Un $\beta$ plus faible rend le \omterm{def:dv:sabr-and-smile-dynamics:backbone}{backbone} plus pentu, ce qui augmente le delta de Hagan pour un call~; $\rho$ recalibre le smile
et modifie le terme de correction~; les deltas se déplacent du \omterm{def:dv:greeks-and-the-hedging-pnl:greeks}{vega} multiplié par la variation de pente.
\textbf{18.} Le \omterm{def:dv:sabr-and-smile-dynamics:mvdelta}{delta de variance minimale} n'est aussi bon que la co-variation de $\alpha$ et de $F$ prédite par le modèle~; il faut
estimer cette co-variation à partir des données et comparer.
\textbf{19.} 19~\% de variance en moins qu'avec le delta de Black (rapport 0.81).
\textbf{20.} Sur la dynamique du modèle~: la façon dont la volatilité évolue avec le sous-jacent.
\end{solution}

\begin{solution}{iq:dv:sabr-and-smile-dynamics:1}
$dF=\alpha F^\beta dW^1$, $d\alpha=\nu\alpha dW^2$, corrélation $\rho$. $\alpha$~: niveau~; $\beta$~: \omterm{def:dv:sabr-and-smile-dynamics:backbone}{backbone} (la façon dont la volatilité
à la monnaie bouge avec $F$)~; $\rho$~: skew autour du \omterm{def:dv:sabr-and-smile-dynamics:backbone}{backbone}~; $\nu$~: courbure des ailes.
\iqlookfor{le rôle de chaque paramètre, et le fait que $\beta$ porte sur la dynamique.}
\end{solution}

\begin{solution}{iq:dv:sabr-and-smile-dynamics:2}
La volatilité à la monnaie en fonction du forward, $\approx\alpha F^{\beta-1}$. On régresse les variations de
$\ln\sigma_{\mathrm{ATM}}$ sur les variations de $\ln F$ sur une période de régime stable~; la pente vaut $\beta-1$.
\iqlookfor{la définition et une procédure d'estimation.}
\end{solution}

\begin{solution}{iq:dv:sabr-and-smile-dynamics:3}
$V(S,v)$~; $dV=V_SdS+V_vdv+\dots$~; on régresse $dv$ sur $dS$~: la pente vaut $\Cov(dv,dS)/\Var(dS)$~; on couvre $V_S+V_v\times$ pente. Le
terme restant n'est pas corrélé à $dS$.
\iqlookfor{l'argument de la régression.}
\end{solution}

\begin{solution}{iq:dv:sabr-and-smile-dynamics:4}
La \omterm{def:dv:local-volatility:model}{volatilité locale} prédit que le smile se déplace en sens inverse du forward, alors que celui du marché des taux évoluait dans le même sens~;
le delta du modèle avait donc une mauvaise correction, et ses couvertures pouvaient être moins bonnes que celles de Black.
\iqlookfor{une dynamique dans le mauvais sens.}
\end{solution}

\begin{solution}{iq:dv:sabr-and-smile-dynamics:5}
La \omterm{def:dv:sabr-and-smile-dynamics:hagan}{formule de Hagan} est un développement asymptotique~; loin de la monnaie et aux longues échéances, elle implique des prix de call non convexes.
On la détecte en calculant la densité, et on évalue dans cette zone avec la densité SABR sans arbitrage (une équation forward unidimensionnelle),
un modèle décalé, ou une extrapolation séparée des ailes.
\iqlookfor{le domaine du développement et une solution.}
\end{solution}

\begin{solution}{iq:dv:sabr-and-smile-dynamics:6}
Les opérations dont la valeur dépend de la façon dont le smile se déplace avec le forward~: options couvertes en delta, options à départ différé
et options de spread entre échéances. On couvre en calculant les risques sous plusieurs couples $(\beta,\rho)$ calibrés sur le même smile
et en gardant une réserve pour l'écart de leurs P\&L, ou en couvrant le \omterm{def:dv:greeks-and-the-hedging-pnl:greeks}{vega} par strike (par compartiments) afin que la dynamique du modèle
pèse moins.
\iqlookfor{produits sensibles à la dynamique et gestion de l'incertitude de modèle.}
\end{solution}
